\documentclass[11pt]{article}

\usepackage[margin=1in]{geometry}
\usepackage{amsmath,amssymb,amsthm}
\usepackage{enumitem}
\usepackage[hidelinks]{hyperref}

\title{Direction 2: a compactness program for convex Schiffer counterexamples}
\author{}
\date{\today}

\newcommand{\R}{\mathbb{R}}
\newcommand{\Om}{\Omega}
\newcommand{\eps}{\varepsilon}
\newcommand{\1}{\mathbf 1}
\newcommand{\nor}{\partial_\nu}
\newcommand{\diam}{\mathrm{diam}}
\newcommand{\sgn}{\mathrm{sgn}}

\newtheorem{lem}{Lemma}
\newtheorem{prop}[lem]{Proposition}
\newtheorem{cor}[lem]{Corollary}
\newtheorem{thm}[lem]{Theorem}
\newtheorem{rem}[lem]{Remark}
\newtheorem{conj}[lem]{Conjecture}
\newtheorem{step}[lem]{Step}

\begin{document}
\maketitle

\section*{Scope}

Direction 2 of \texttt{ideas-schiffer.tex} is the \emph{unweighted}
compactness program for the planar Schiffer conjecture. The previous
notes in this folder (\texttt{direction2-barrel-limit},
\texttt{direction2-heat-extension-computation},
\texttt{direction2-parabolic-to-classical},
\texttt{direction2-stationary-weighted-perturbation}) explored a
log-concave extension of the problem; that line eventually showed the
parabolic angle is too weak to close the conjecture
(Theorem~1 of \texttt{direction2-parabolic-to-classical} reduces it
to the classical statement) and is therefore a detour relative to
Direction 2 as written. This note restarts cleanly: it sets up the
compactness program for genuine planar convex Schiffer
counterexamples and identifies precisely where the tools from
arXiv:2412.06344 enter as enabling machinery.

\section{The compactness program in precise form}

\begin{conj}[Schiffer for planar convex domains]\label{conj:schiffer}
If $\Om\subset\R^2$ is a bounded smooth strictly convex domain
admitting a nonzero solution of
\begin{equation}\label{eq:schiffer}
    \Delta u+\lambda u=0\ \text{in }\Om,\qquad
    \nor u=0\ \text{on }\partial\Om,\qquad
    u|_{\partial\Om}=c\in\R,
\end{equation}
for some $\lambda>0$, then $\Om$ is a disc.
\end{conj}

The compactness strategy of Direction 2: assume the negation, take a
sequence $\{(\Om_j,u_j,\lambda_j)\}$ of normalized non-disc convex
counterexamples, extract a limit by Blaschke selection, pass the
overdetermined data to the limit, and derive a contradiction. The
program splits cleanly into four steps.

\begin{step}[Setup and normalization]\label{step:normalize}
Choose normalizations on the sequence so that Blaschke selection
gives a non-degenerate convex limit $\Om_\infty$, and so that the
eigenvalues $\lambda_j$ are bounded in a compact subinterval of
$(0,\infty)$.
\end{step}

\begin{step}[Spectral compactness]\label{step:spectral}
Show that the normalized eigenfunctions $u_j$, after possibly
extracting subsequences, converge in a strong enough topology
(uniformly with bounded derivatives) to a limit $u_\infty$ on
$\Om_\infty$, and that $\lambda_j\to\lambda_\infty>0$ with
$u_\infty\not\equiv 0$.
\end{step}

\begin{step}[Passing the overdetermination]\label{step:pass}
Show that $u_\infty$ satisfies \eqref{eq:schiffer} on $\Om_\infty$,
including both overdetermined conditions
$\nor u_\infty=0$ and $u_\infty|_{\partial\Om_\infty}=c_\infty$.
\end{step}

\begin{step}[Disposing of the disc limit]\label{step:disc}
Either rule out $\Om_\infty=$ disc by a local rigidity theorem near
the disc, or rule out a non-disc $\Om_\infty$ by an extremality
argument applied to a minimizing sequence.
\end{step}

\section{Step 1: normalizations}

Three families of normalization are compatible with Schiffer's
overdetermined problem, since \eqref{eq:schiffer} is invariant under
$(\Om,u,\lambda)\mapsto(\tau\Om,\tau\circ u,\lambda)$ for rigid motions
$\tau$ and under $(\Om,u,\lambda)\mapsto(r\Om,u(\cdot/r),\lambda/r^2)$
for dilations. Centre at $0$ to remove translations; rotate so that
the diameter axis is horizontal to remove rotations. To remove the
dilation freedom, fix one scalar geometric quantity.

\begin{itemize}[leftmargin=*]
    \item \textbf{Diameter normalization}: $\diam(\Om_j)=1$. Then
    $\Om_j\subset B_1(0)$ and Blaschke gives a convex Hausdorff limit
    $\Om_\infty$. The drawback is that $\Om_j$ can degenerate to a
    line segment (zero width).
    \item \textbf{Area normalization}: $|\Om_j|=\pi$ (the area of
    a unit disc). Now $\Om_j$ can have arbitrarily large diameter
    (long thin rectangle), so an additional bound is required to
    apply Blaschke.
    \item \textbf{Combined normalization}: $|\Om_j|=\pi$ \emph{and}
    $\diam(\Om_j)\le D_*$, with $D_*$ a constant chosen below.
\end{itemize}

\begin{lem}[Diameter bound from a fixed eigenvalue]\label{lem:diam}
If $\lambda_j\le\Lambda$ for all $j$ and $|\Om_j|=\pi$, then there is
$D_*=D_*(\Lambda)$ such that $\diam(\Om_j)\le D_*$.
\end{lem}

\begin{proof}
The function $u_j$ is a Neumann eigenfunction of $-\Delta$ at
eigenvalue $\lambda_j$ on the convex $\Om_j$. The Payne--Weinberger
inequality gives $\lambda_j\ge\pi^2/\diam(\Om_j)^2$ for the first
nonzero Neumann eigenvalue, but $u_j$ is not necessarily a
\emph{first} eigenfunction. However, the number of Neumann
eigenvalues below $\Lambda$ on $\Om_j$ is at most $C(\Lambda)|\Om_j|=
C(\Lambda)\pi$ by Weyl's asymptotics (uniform under area
normalization for convex domains by Li--Yau bounds, see below). So
$\lambda_j$ is one of at most $C(\Lambda)\pi$ eigenvalues, and the
smallest such eigenvalue dominates the others up to a multiplicative
constant. Combined with Payne--Weinberger applied to the first
eigenvalue, $\diam(\Om_j)$ is bounded by $D_*(\Lambda)$.
\end{proof}

Combining the area normalization with an a priori bound on
$\lambda_j$ produces a compact family in the Hausdorff topology,
modulo the question whether $\lambda_j$ is bounded.

\begin{rem}[Upper eigenvalue bound]\label{rem:upperbound}
Schiffer's overdetermined problem is unrestricted in $\lambda$, so
without further input $\lambda_j$ is not bounded. Three natural ways
to obtain $\lambda_j\le\Lambda$:
\begin{itemize}[leftmargin=*]
    \item Restrict attention to counterexamples with $\lambda$ below
    a chosen threshold $\Lambda$ (e.g.\ to the first $K$ eigenvalues).
    Deng's results already cover $K=1,\ldots,k_0$ for small $k_0$ on
    convex planar domains.
    \item Minimize $\lambda$ over the class of counterexamples and
    bound the infimum by a comparison construction. The infimum is
    attained at a limit point of a minimizing sequence; the program
    naturally takes this limit.
    \item Bound $\lambda_j$ in terms of the area-normalized
    fundamental Neumann gap by an isoperimetric inequality (a
    Fauconnier-type bound is conjectural here).
\end{itemize}
The second option is the cleanest extremality choice and is the one
adopted below.
\end{rem}

\begin{conj}[Schiffer infimum]\label{conj:inf}
Define
\[
    \Lambda_*:=\inf\{\lambda\in(0,\infty):\,
        \exists\Om\subset\R^2\text{ convex, non-disc, smooth, area }\pi,
        \text{ with a nonzero }u\text{ solving \eqref{eq:schiffer}}\}.
\]
Then $\Lambda_*<\infty$.
\end{conj}

A simple comparison upper bound suffices: any ellipse close to a disc
should produce \emph{near-Schiffer} data with small Schiffer defect,
which a perturbative bifurcation argument would either complete to a
true Schiffer pair or rule out (Direction 3 territory). Assuming
$\Lambda_*<\infty$, fix a minimizing sequence
$\{(\Om_j,u_j,\lambda_j)\}$ with $|\Om_j|=\pi$,
$\lambda_j\to\Lambda_*$, and apply Lemma~\ref{lem:diam} with
$\Lambda=\Lambda_*+1$ to obtain $\diam(\Om_j)\le D_*$.

\section{Step 2: spectral compactness via the 2412 toolkit}

This is the step where arXiv:2412.06344 contributes most directly.
de Dios Pont's Step 4 (after Theorem 2.6) is: \emph{``By Li--Yau's
uniform Harnack inequality, the eigenfunctions should be
equicontinuous, and thus have convergent subsequences in $C^0$.''}
The cited references are [LY86] (Li--Yau, the parabolic kernel of
the Schr\"odinger operator) and [Wan06] (Wang, gradient estimates).
The relevant fact for us is:

\begin{thm}[Uniform gradient estimate, Wang]\label{thm:wang}
Let $\Om\subset\R^n$ be a smooth bounded convex domain and let $u$
solve $-\Delta u=\lambda u$ in $\Om$ with $\nor u=0$ on $\partial\Om$.
Then
\begin{equation}\label{eq:wang}
    \sup_\Om|\nabla u|^2\le C_n\,\lambda\,\sup_\Om u^2,
\end{equation}
with $C_n$ depending only on the dimension $n$.
\end{thm}

The constant in \eqref{eq:wang} does \emph{not} depend on the
specific shape of $\Om$, only on dimension. This is the
``uniform Harnack across convex domains'' that the 2412 paper uses.

\begin{cor}[Equicontinuity of normalized Schiffer eigenfunctions]\label{cor:equi}
Let $\{(\Om_j,u_j,\lambda_j)\}$ be a sequence of smooth convex
Schiffer counterexamples in $\R^2$ with $|\Om_j|=\pi$,
$\diam(\Om_j)\le D_*$, and $\lambda_j\le\Lambda_*+1$. Normalize
$\|u_j\|_{L^\infty(\Om_j)}=1$. Then $\{u_j\}$ is equicontinuous and
uniformly bounded on the compact set $B_{D_*}(0)\supset\Om_j$ (after
zero-extension or a Whitney extension off $\Om_j$).
\end{cor}

\begin{proof}
$\sup|\nabla u_j|\le \sqrt{C_2\lambda_j}\le \sqrt{C_2(\Lambda_*+1)}$
by Theorem~\ref{thm:wang} applied with $n=2$. The $u_j$ are
$C^{0,1}$-bounded with constant $\sqrt{C_2(\Lambda_*+1)}$
independent of $j$. Extend $u_j$ by zero off $\Om_j$ to obtain a
sequence of $C^{0,1}$-bounded functions on the compact set
$\overline{B_{D_*}(0)}$, hence equicontinuous and uniformly bounded.
\end{proof}

\begin{cor}[Spectral subsequential limit]\label{cor:speclim}
Under the hypotheses of Corollary~\ref{cor:equi}, after extraction:
\[
    \Om_j\to\Om_\infty\text{ in Hausdorff distance},\qquad
    u_j\to u_\infty\text{ uniformly on compacta},\qquad
    \lambda_j\to\lambda_\infty.
\]
$\Om_\infty$ is convex with $|\Om_\infty|\le\pi$ and
$\diam(\Om_\infty)\le D_*$. If $\Om_\infty$ has nonempty interior,
then $u_\infty\in C^{0,1}(\Om_\infty)$ solves
$-\Delta u_\infty=\lambda_\infty u_\infty$ in $\Om_\infty^\circ$.
\end{cor}

\begin{proof}
Arzel\`a--Ascoli on $\{u_j\}$, Blaschke on $\{\Om_j\}$, and bounded
$\{\lambda_j\}$ all give subsequential limits. The PDE passes
to the limit in distributions on every compact subset of
$\Om_\infty^\circ$ because of the uniform Lipschitz bound and
uniform $\lambda$ bound.
\end{proof}

The remaining question is whether $|\Om_\infty|>0$, equivalently
whether $\Om_\infty$ is not collapsed to a segment. This is the
\emph{non-degeneracy} of the limit and is the main technical issue.

\begin{lem}[Non-degeneracy from a width lower bound]\label{lem:nondeg}
If the inradius $r(\Om_j)\ge r_*>0$ uniformly, then $\Om_\infty$ has
nonempty interior.
\end{lem}

The hypothesis $r(\Om_j)\ge r_*$ requires either a structural
argument (e.g.\ thin domains have large Schiffer eigenvalue,
contradicting $\lambda_j\le\Lambda_*+1$) or an explicit further
normalization. The cleanest sufficient condition:

\begin{conj}[Quantitative thin-strip exclusion]\label{conj:thinstrip}
There is a function $\omega:(0,\infty)\to(0,\infty)$ with
$\omega(r)\to\infty$ as $r\to 0^+$ such that for every smooth convex
Schiffer counterexample $\Om$ in $\R^2$ with $|\Om|=\pi$ and
$r(\Om)\le r$, every Schiffer eigenvalue $\lambda$ of $\Om$ satisfies
$\lambda\ge\omega(r)$.
\end{conj}

This is plausible by analogy with Faber--Krahn on thin domains:
Dirichlet eigenvalues on a thin strip blow up like (width)$^{-2}$. A
Schiffer counterexample by definition has $u|_{\partial\Om}=c$, and
unless $c=0$ this means the function does not vanish at the boundary,
so the Dirichlet bound does not literally apply. A correct
substitute is the Robin-type Faber--Krahn or a careful estimate on
$u-c$, which solves $\Delta(u-c)+\lambda u=0$, $\nor(u-c)=0$,
$(u-c)|_{\partial\Om}=0$. This is a mixed problem whose eigenvalues
admit thin-strip lower bounds.

Conditional on Conjecture~\ref{conj:thinstrip}, the limit
$\Om_\infty$ has positive inradius and Step 2 closes.

\section{Step 3: passing the overdetermined conditions}

With $u_j\to u_\infty$ uniformly and $\Om_j\to\Om_\infty$ in
Hausdorff distance, the Dirichlet trace is easy:
$u_j|_{\partial\Om_j}=c_j$. The constants $c_j$ are bounded
($|c_j|\le\sup|u_j|=1$) and pass to $c_\infty\in[-1,1]$. By uniform
convergence and Hausdorff continuity of the boundary,
$u_\infty|_{\partial\Om_\infty}=c_\infty$.

The Neumann trace is harder. $\nor u_j=0$ on $\partial\Om_j$, and we
want $\nor u_\infty=0$ on $\partial\Om_\infty$. This requires
$C^1$-convergence of $u_j$ near the boundary and Hausdorff
convergence of the normal vector fields, which fails at corners of
$\Om_\infty$.

If $\partial\Om_\infty$ is $C^1$ (which follows from strict convexity
+ Hausdorff convergence of strictly convex bodies on compact subsets
of the boundary), then $\nor u_\infty=0$ follows from interior
gradient estimates and boundary $C^{1,\alpha}$ regularity for elliptic
equations on convex domains (Krylov--Safonov / Caffarelli).

The two cases that fail are:
\begin{itemize}[leftmargin=*]
    \item $\partial\Om_\infty$ has corners. Then $\Om_\infty$ is a
    convex polygon (or has line-segment pieces), $\nor u_\infty$ is
    well-defined only away from the singularities, and the
    overdetermined condition holds in a weak sense.
    \item $\Om_\infty$ has flat boundary pieces. Then $\partial\Om_\infty$
    contains line segments and the limit is no longer strictly
    convex.
\end{itemize}
Both can be handled by treating the singular part of $\partial\Om_\infty$
as exceptional.

\begin{prop}[Limit Schiffer pair, conditional]\label{prop:limitpair}
Conditional on Conjecture~\ref{conj:thinstrip} and the absence of
corners in $\partial\Om_\infty$, the limit
$(\Om_\infty,u_\infty,\lambda_\infty)$ is a Schiffer pair: $\Om_\infty$
is a smooth convex domain, $\lambda_\infty\ge\Lambda_*$, and
$u_\infty$ is a nontrivial Neumann eigenfunction of $-\Delta$ on
$\Om_\infty$ satisfying $u_\infty|_{\partial\Om_\infty}=c_\infty$.
\end{prop}

By construction $\lambda_\infty=\Lambda_*$ (since the sequence is
minimizing). $\Om_\infty$ is therefore a Schiffer-minimizer: a convex
domain of area $\pi$ at which the Schiffer infimum is attained.

\section{Step 4: dichotomy at the minimizer}

There are two cases on the limit.

\paragraph{Case A: $\Om_\infty=$ disc.} The unit disc satisfies
Schiffer at each radial Neumann eigenvalue
$\lambda_k^{\mathrm{rad}}=(j'_{0,k})^2$. So $\Lambda_*\ge
(j'_{0,1})^2=(j_{1,1})^2$, and equality would mean $\Om_\infty=B$ at
the first radial Neumann eigenvalue. The sequence $\Om_j\to B$ in
Hausdorff distance with $\lambda_j\to(j_{1,1})^2$ and
$u_j\to J_0(j_{1,1}r)$ (radial Bessel mode).

This case is excluded by a \emph{local rigidity theorem near the
disc}: there exists $\delta>0$ such that no smooth convex
non-disc $\Om$ with $d_H(\Om,B)<\delta$ admits a Schiffer
eigenfunction. Existing perturbative work (Canuto, Kobayashi) gives
such results for low Neumann eigenvalues. Whether the existing
estimates cover the radial Schiffer eigenvalues
$\lambda_k^{\mathrm{rad}}$ specifically is the technical question. A
clean version of Direction 3 (bifurcation from the disc) would
either confirm this or provide an explicit non-disc branch.

\paragraph{Case B: $\Om_\infty\ne$ disc.} Then $\Om_\infty$ is a
smooth convex non-disc Schiffer counterexample of area $\pi$ at
eigenvalue $\Lambda_*$. By minimality, $\Om_\infty$ satisfies an
Euler--Lagrange-type condition under area-preserving convex
deformations. Specifically, every smooth one-parameter family of
area-preserving deformations $\Om_s$ with $\Om_0=\Om_\infty$ that
preserves Schiffer (and hence has $\lambda(\Om_s)$ well-defined as a
Schiffer eigenvalue) must satisfy
\(
    \partial_s\lambda(\Om_s)\big|_{s=0}\ge 0.
\)

The shape-derivative formula on the disc-type domains is
\[
    \partial_s\lambda=-\int_{\partial\Om}|\nabla u|^2\,V_n\,dS,
\]
where $V_n$ is the normal velocity. On a Schiffer counterexample
$|\nabla u|^2$ on $\partial\Om$ is determined by the Cauchy data:
$\nor u=0$ and $u=c$ together with $-\Delta u=\lambda u$ at the
boundary give
$|\nabla u|^2|_{\partial\Om}=(\partial_T u)^2=0$ (since $u|_{\partial\Om}=c$
forces zero tangential derivative). Combined with $\nor u=0$, the
full gradient vanishes on $\partial\Om$ and the shape derivative
vanishes identically. So minimality in $\lambda$ gives no information
at first order.

At \emph{second} order, the variational inequality reads (for a
Schiffer-preserving deformation $V_n=f$ on $\partial\Om$)
\[
    \partial_s^2\lambda\big|_{s=0}\ge 0,
\]
and the second shape derivative for a Schiffer pair has been computed
(in different language) by Garofalo--Lewis, Sirakov, and others. The
resulting nonnegativity is the right structural constraint, and it
ought to be incompatible with $\Om_\infty\ne$ disc when combined with
the spectral pair structure.

\section{What 2412 actually gives us, and what is still needed}

\begin{enumerate}[leftmargin=*]
    \item \textbf{Used directly from 2412.}
    \begin{itemize}[leftmargin=*]
        \item Theorem~\ref{thm:wang} (Wang/Li--Yau uniform gradient
        estimate). This is the key tool for Step 2.
        Equicontinuity of Schiffer eigenfunctions across the
        sequence is the precise analog of de Dios Pont's Step 4.
        \item Smooth convex approximation (Lemma A.2 of 2412). This
        lets us replace non-smooth limits $\Om_\infty$ by smooth
        approximations at every step without loss.
    \end{itemize}

    \item \textbf{Borrowable from 2412 but not yet imported.}
    \begin{itemize}[leftmargin=*]
        \item The Brunn--Minkowski / log-concave inequalities used in
        2412 to control the eigenvalue $\lambda_{\Om,V}$ via
        weighted Faber--Krahn. The unweighted version of these
        inequalities is exactly what Conjecture~\ref{conj:thinstrip}
        needs: thin convex domains have a quantitative Schiffer
        eigenvalue lower bound.
        \item The cap/lateral split (Section~2 of 2412). For
        Direction 2 this translates to a dichotomy between
        ``elongated'' and ``compact'' limits: an elongated
        $\Om_\infty$ corresponds to the lateral collapse, a thin one
        to the cap collapse. The 2412 analysis says lateral
        collapse is the rigidity-friendly case.
    \end{itemize}

    \item \textbf{Not in 2412, still needed.}
    \begin{itemize}[leftmargin=*]
        \item Local rigidity near the disc (Case A of Step 4).
        This belongs to Direction 3.
        \item The second-variation positivity at a minimizer
        (Case B). Needs an explicit Hessian-of-shape calculation
        for Schiffer pairs.
        \item Existence of a minimizing sequence with bounded
        eigenvalue (Conjecture~\ref{conj:inf}). This requires
        constructing a near-Schiffer family with bounded $\lambda$,
        which is essentially a perturbative existence statement
        (Direction 3 again).
    \end{itemize}
\end{enumerate}

\section{Concrete next deliverables}

\begin{enumerate}[leftmargin=*]
    \item \textbf{Quantitative thin-strip exclusion} (Conjecture
    \ref{conj:thinstrip}). Prove that thin convex domains have large
    Schiffer eigenvalues, by reducing to a Dirichlet eigenvalue
    estimate for $u-c$ via a Robin-type Faber--Krahn. This is a
    self-contained technical lemma; one page if it works.
    \item \textbf{Apply Theorem~\ref{thm:wang} explicitly.}
    Write down the uniform Lipschitz bound for the Schiffer
    eigenfunction $u$ on a convex domain in terms of
    $\lambda$ and $|\Om|$, and use it to verify equicontinuity in
    Corollary~\ref{cor:equi}.
    \item \textbf{Second shape derivative on a Schiffer pair.}
    Compute the Hessian of $\lambda$ as a function of the boundary
    deformation, restricted to Schiffer-preserving deformations
    around a putative non-disc minimizer. Either show that the
    Hessian is sign-definite (ruling out Case B) or identify an
    explicit Schiffer-preserving direction that decreases $\lambda$,
    contradicting minimality.
    \item \textbf{Verify the existing local rigidity near the disc.}
    Survey Canuto, Kobayashi, and any subsequent work to determine at
    which Neumann eigenvalues local rigidity is currently available,
    and at which eigenvalues it remains open. The radial Neumann
    eigenvalues $(j'_{0,k})^2$ are the relevant target.
\end{enumerate}

\end{document}
